\section{Thin-wall phase matching and a planar transmission zero}\label{sec:thinwall}
Put $\epsilon=\omega^2-\omega_c^2$ and $S_c=1/(2\beta)$.
The leading thin-wall profile and radius in three dimensions are
\begin{equation}\label{eq:thinwall}
R_{\mathrm{tw}}\simeq\frac{1}{2\sqrt\beta\,\epsilon},\qquad
f^2(r)\simeq\frac{S_c}{1+\exp[(r-R_{\mathrm{tw}})/\sqrt\beta]}.
\end{equation}
These are asymptotic approximations, not substitutes for the certified
profiles. The leading radius relation is established prior work: with
$\Delta=m_\phi^2-\omega_0^2$, Heeck et al. use
$\kappa^2=(\omega^2-\omega_0^2)/\Delta$ and
$R^*=\sqrt\Delta R$, and obtain $R^*\simeq1/\kappa^2$
in their Eq.~(40)~\cite{Heeck2021}.
Here $\Delta=1/(4\beta)$, which gives the radius in
Eq.~\eqref{eq:thinwall}. Their inflection-point radius and our
half-density radius need not have identical subleading coefficients.
Their Eq.~(58) gives
$R^*=\kappa^{-2}+\tfrac14-\tfrac{5}{16}\kappa^2+O(\kappa^4)$;
these terms require a radius-convention conversion before comparison
with the present half-density radius~\cite{Heeck2021}.
Diagonalizing the constant interior channel matrix gives the
propagating interior wavenumber
\begin{equation}\label{eq:interior}
k_{\mathrm{in}}^2=\omega^2+\rho^2-D_c+
 \sqrt{4\omega^2\rho^2+C_c^2},\qquad
D_c=1+\frac1{4\beta},\quad C_c=\frac1{2\beta}.
\end{equation}
At the stated interior amplitude, this is the same dispersion branch as
Eq.~(101) of Zhang, Zhou, and Zhu~\cite{ZhangZhouZhu2025}; the interior
wavenumber is an established ingredient, not a new result here.
A regular interior standing wave reaches the wall with phase
$\Psi=k_{\mathrm{in}}R_{\mathrm{tw}}$. The bare closed-channel wall approximation
supplies a localized state, while matching it to the open interior
wave permits a sign-changing radiation amplitude. A useful leading
parametrization is
\begin{equation}\label{eq:phase}
k_{\mathrm{in}}(\omega_n,\rho_n)R_{\mathrm{tw}}(\omega_n)
 \simeq\pi[n+\theta(\epsilon_n)],\qquad
\rho_n\simeq\omega_n+c(\epsilon_n).
\end{equation}
The radius and interior dispersion follow from the model approximations;
the wall phase $\theta$ and dressed level shift contained in $c$ remain
calibrated. A separate coupled planar-wall calculation now supplies a
numerical transmission zero, but does not determine both finite-curvature
functions independently. Consequently, agreement of the calibrated
phase rule with the sequence is not a parameter-free prediction.

For $\beta=1/2$ and $\omega^2=1/2$, the exact planar profile is
$S(x)=[1+\exp(\sqrt2 x)]^{-1}$. Solving the two-channel scattering
problem, with one open and one evanescent channel on each side, gives
\begin{equation}\label{eq:planarzero}
\rho_z\simeq1.52414976213,\qquad
k_{\mathrm{in}}(\rho_z)\simeq1.92332459129.
\end{equation}
An independently implemented Numerov boundary-value calculation confirmed
this zero before reading the comparison calculation. It used a continuously
phase-normalized transmission amplitude, actual complex-transmission
residuals, flux and reciprocity controls, and separate mesh and boundary
refinements. A conservative empirical frequency uncertainty is $10^{-9}$;
this is not a certified enclosure. Two finite frequency scans found no
additional zero, but do not prove uniqueness throughout the open-channel
window. The planar scattering state has an incident and reflected wave;
it is not itself a square-integrable embedded mode.

If $k_{\mathrm{in}}\to k_\infty>0$ and the wall phase has a finite limit,
Equation~\eqref{eq:phase} predicts, within the approximation,
\begin{equation}\label{eq:spacing}
\epsilon_n\sim\frac1{b_\infty n},\qquad
b_\infty=\frac{2\sqrt\beta\,\pi}{k_\infty}.
\end{equation}
From the historical rounded locations $0.535449$, $0.532772$,
$0.530470$, $0.528469$, the inverse-detuning steps for indices 12--15
are approximately 2.3043, 2.3053 and 2.3068;
these descriptive values do not specify accuracy beyond the input locations.
The matching locations added in this revision (Section~\ref{sec:sequence})
give $2.3045$, $2.3053$ and $2.3059$ for the same three steps.
The calibrated model gives 2.298, while its bare
wall-state estimate gives 2.334. Substituting the independently computed
planar zero into Equation~\eqref{eq:spacing} instead gives
$b_\infty^{\mathrm{wall}}\simeq2.31000162857$, with a conservative empirical
numerical uncertainty of $2\times10^{-9}$. This wall-based candidate was not
fitted to ladder locations. Identifying it with the actual radial limiting
step still requires a controlled curvature and phase-matching limit.
This supports a useful spacing description;
it does not prove existence of infinitely many localized modes or identify
an exact limiting constant. The phase model also motivates alternating
signs of a matching Jacobian and quadratic linewidth minima, but neither
property independently replaces a tail-controlled existence proof.

The open part of this description can be stated precisely.
\begin{conjecture}[radial ladder and limiting spacing]\label{conj:ladder}
Let $\beta=1/2$ and $\ell=0$. There is an index $n_0$ such that for
every $n\ge n_0$ there is a parameter pair $(\omega_n,\rho_n)$ with
$\omega_c<\omega_n<1$ and $1-\omega_n<\rho_n<1+\omega_n$ at which
Eq.~\eqref{eq:proof-profile} has a positive global solution with
$f(r)\to0$ and Eq.~\eqref{eq:proof-radial-system} has a nonzero real
solution regular at the origin that obeys the decay
bound~\eqref{eq:proof-decay} for some $K_J$ and $L$, which may depend
on $n$. These pairs can be chosen such that $\epsilon_n=\omega_n^2-1/2$
decreases strictly, $\epsilon_n\to0$, $\rho_n\to\rho_z$, and
\begin{equation*}
\epsilon_{n+1}^{-1}-\epsilon_n^{-1}\longrightarrow
b_\infty^{\mathrm{wall}}=\frac{2\sqrt\beta\,\pi}{k_{\mathrm{in}}(\rho_z)}
\simeq2.3100 .
\end{equation*}
Here $\rho_z$ denotes the conjectured simple zero of the exact planar
transmission amplitude near the numerical value of
Eq.~\eqref{eq:planarzero}, and $k_{\mathrm{in}}$ is evaluated at
$\omega=\omega_c$.
\end{conjecture}
The evidence for Conjecture~\ref{conj:ladder} is numerical and
heuristic: certificates for $n=1,2$; finite-domain candidates through
$n=18$ whose inverse-detuning steps rise slowly, from $2.3036$ to
$2.3059$ between indices $11$ and $15$ in the matching locations of
Section~\ref{sec:sequence}; and the uncertified planar transmission
zero. Three ingredients of a proof are missing: (i)~an enclosure of the
planar zero $\rho_z$ together with a nondegeneracy (transversality)
condition; (ii)~estimates, uniform as $\epsilon\to0$, that control the
difference between the rescaled curved-wall matching map and its planar
limit, including the deviation of the profile from the planar kink and
the $1/r$ and centrifugal terms; and (iii)~a quantization estimate
controlling successive phase increments, sufficient to establish
convergence of the inverse-detuning gaps, together with a degree or
fixed-point argument that converts (i) and (ii) into one localized mode
per index.
None of them is attempted here. The finite steps are compatible with
convergence to $2.3100$ but do not demonstrate it; for $\beta=1$ the
reported steps ($2.5964$ to $2.6038$, below) remain about $0.6$--$0.9\,\%$
below their planar candidate $2.6186$.
The inverse-first-power scaling follows from the calibrated model, not
from an independently established asymptotic exponent. Exploratory free
power fits yielded exponents between $-0.76$ and $-0.88$; the reciprocal
fit improved the comparison metric only by factors of approximately
$1.2$--$2$, below the originally sought factor of ten. Finite-data agreement
therefore does not establish an exact asymptotic law.

A local numerical check also supports a planar transmission cancellation
for the modified potential $U(S)=S-S^2+S^3$ at
$\omega_{\min}=\sqrt{3}/2$. An independent Numerov implementation,
with the candidate frequency already known, gives
$\rho_z\simeq1.77345307180$ and $k_{\mathrm{in}}\simeq2.3994321234$.
The last grid change, $h=0.04$ to $0.02$, shifts $\rho_z$ by
$1.64\times10^{-10}$, while increasing the half-domain from $32$ to $40$
at the finest spacing changes it by $4.2\times10^{-14}$; these are
empirical sensitivities, not certified bounds. The corresponding thin-wall
spacing candidate is $2\pi/k_{\mathrm{in}}\simeq2.618613482$. This known-target
local replication neither counts all planar zeros nor establishes the
radial ladder limit for the modified potential.

The corresponding radial calculation is also available. Its last
previously reported inverse-detuning step was $2.596420$, about $0.9\%$
below the planar candidate. A subsequent search gave the following
finite-domain numerical locations; $k$ continues that study's local
index and is not a certified node count:
\begin{center}\small
\begin{tabular}{rlll}\toprule
$k$&$\omega^2$&$1/(\omega^2-3/4)$&Resolved winding\\\midrule
0&$0.778585440$&$34.98284$&$-1$\\
$-1$&$0.776606628$&$37.58462$&$+1$\\
$-2$&$0.774882776$&$40.18844$&$-1$\\\bottomrule
\end{tabular}\end{center}
The reported empirical uncertainties are about $2\times10^{-8}$ in
$\omega^2$ and $3\times10^{-5}$ in inverse detuning. These later
increments are approximately $2.6018$ and $2.6038$; convergence to the
planar value remains a conjectural interpretation. A frozen
absolute-location formula truncated after the $O(1)$ term failed its
combined test: inverse-detuning errors for this potential were
$0.23$--$0.26$, outside the preassigned $0.10$ window; an extrapolation
of the preceding radial locations was more accurate for both potentials.
This failure concerns the finite-order absolute-location formula,
not a proof or disproof of the leading wall spacing.
